\documentclass[11pt,a4paper]{article}

\usepackage[utf8]{inputenc}
\usepackage[T1]{fontenc}
\usepackage[margin=2cm]{geometry}
\usepackage{parskip}
\usepackage{titlesec}
\usepackage{enumitem}
\usepackage{booktabs}
\usepackage{xcolor}
\usepackage{tcolorbox}
\usepackage{graphicx}
\usepackage{hyperref}
\usepackage{fancyhdr}
\usepackage{listings}
\usepackage{longtable}
\usepackage{array}

\tcbuselibrary{skins,breakable}

\definecolor{accent}{RGB}{30,80,150}
\definecolor{softgray}{RGB}{245,246,248}
\definecolor{codebg}{RGB}{248,248,250}
\definecolor{warn}{RGB}{160,60,40}

\titleformat{\section}{\Large\bfseries\color{accent}}{\thesection}{0.8em}{}
\titleformat{\subsection}{\large\bfseries}{\thesubsection}{0.6em}{}

\hypersetup{colorlinks=true, linkcolor=accent, urlcolor=accent,
  pdftitle={AgentIA --- The Spec-Driven Design},
  pdfauthor={AgentIA Engineering}}

\pagestyle{fancy}
\fancyhf{}
\fancyhead[L]{\small\color{accent}AgentIA --- Spec-Driven Design}
\fancyhead[R]{\small\color{accent}\thepage}
\renewcommand{\headrulewidth}{0.4pt}

\lstset{
  basicstyle=\ttfamily\footnotesize, backgroundcolor=\color{codebg},
  frame=single, rulecolor=\color{softgray}, breaklines=true,
  columns=fullflexible, keepspaces=true, showstringspaces=false,
  xleftmargin=0.4em, xrightmargin=0.4em,
}

\newtcolorbox{notebox}[1][]{colback=softgray,colframe=accent,boxrule=0.4pt,
  arc=2pt,left=6pt,right=6pt,top=4pt,bottom=4pt,breakable,#1}

\title{\vspace{-1.4cm}\textbf{AgentIA --- The Spec-Driven Design}\\[0.25cm]
\large Every specification the platform reads, and where each one is consumed}
\author{AgentIA Engineering}
\date{2026-09-29}

\begin{document}
\maketitle
\thispagestyle{fancy}

\begin{notebox}
\textbf{What this document is.} The platform is spec-first: it ingests formal
specifications, generates against them, and verifies against them. This document
catalogues every specification in force --- the blueprint, the Spec-Kit document,
the constitution, the per-stage instruction documents, the dependency allowlist, the
contract suite, and the skill format --- and states where each is read. The design
claim is that the specification, not the model, is the authority.
\end{notebox}

\tableofcontents
\vspace{0.4cm}

\section{The spine}

Generation is a function of a specification: \emph{spec $\rightarrow$ generate
$\rightarrow$ verify against the spec}. The model proposes; the specification
constrains. The specifications below are the parts of that function, in order of how
directly they determine the output.

\section{The blueprint (the primary spec)}

The \texttt{ArchitectureBlueprint} (\texttt{models/blueprint.py}) is the machine-checkable
specification every generation session starts from. Its schema:

\begin{center}
\begin{longtable}{@{}p{0.34\textwidth}p{0.60\textwidth}@{}}
\toprule
\textbf{Field} & \textbf{Meaning} \\
\midrule
\texttt{serviceName} & lowercase-hyphen service identity \\
\texttt{packageName} & base Java package \\
\texttt{basePort} & the port the service binds (1024--65535) \\
\texttt{databaseMode} & datasource dialect \\
\texttt{layers[]} & the prescribed layer list (controller/service/repository/model) \\
\texttt{entities[]} & the domain entities \\
\texttt{userStories[]} & the acceptance stories (Given/When/Then) \\
\bottomrule
\end{longtable}
\end{center}

Each entity carries \texttt{name}, \texttt{tableName}, and \texttt{attributes}
(name, type, nullable, isPrimaryKey, validationRules); each story carries \texttt{id},
\texttt{priority}, \texttt{role}, \texttt{intent}, \texttt{benefit}, and
\texttt{scenarios} (scenarioId, given, when, then). \texttt{validate\_blueprint}
enforces structural and constitutional integrity before a session may start.

\section{The Spec-Kit document (the human-facing spec)}

A human writes a \texttt{spec.md} conforming to Spec Kit; \texttt{spec\_service.parse\_spec\_markdown}
parses it into an \texttt{ArchitectureBlueprint}, \texttt{validate\_blueprint} checks
it, and \texttt{save\_specification} stores it under a \texttt{specId}. This is the
button visible in the UI as \emph{``Fase 4: Ingesta y Validación Formal de Blueprints
(Spec Kit)''}, and the export button in the Requirements tab
(\emph{``Descargar documento spec.md compatible con Spec Kit''}).

\section{The constitution (the governance spec)}

\texttt{.specify/memory/constitution.md} (v1.1.0) is the six-principle authority the
generator and the compliance gate both defer to:

\begin{enumerate}[leftmargin=1.4em]
  \item \textbf{Strict layering} --- controller $\rightarrow$ service $\rightarrow$
  repository $\rightarrow$ model, each layer talking only to the one below it.
  \item \textbf{Immutable contracts, early validation} --- records only, no JPA
  entities at the boundary, Jakarta Validation on requests.
  \item \textbf{Centralised errors} --- one \texttt{@RestControllerAdvice}, one uniform
  RFC 7807 shape.
  \item \textbf{Offline determinism} --- hermetic \texttt{mvn test -o} in a sealed
  container; no undeclared dependencies, no build-time downloads.
  \item \textbf{Quality gates and bounded repair} --- 100\% test pass; at most three
  auto-repair iterations, then \texttt{BLOCKED}.
  \item \textbf{Secrets and the LangGraph boundary} --- no hardcoded keys; the LLM
  orchestrator stays outside the generated service; no LLM calls from tests.
\end{enumerate}

\section{The instruction documents (the procedural spec)}

The five stage documents (\texttt{resources/instructions/*.md}) are a per-stage
procedural specification: each prescribes the technology contract, the rules, the
output contract, and the prohibitions for one stage. The manifest maps stage to
document; the set is versioned content-addressed, so an artifact records exactly
which instruction text produced it. (Full text in the companion report
\emph{Agents, Prompts and the LangGraph}.)

\section{The dependency allowlist (the build spec)}

\texttt{resources/dependency\_allowlist.json} is the build's specification: the parent,
the five allowed dependencies, and the three build plugins. Anything not listed is a
blocking violation at the scaffold stage, because the offline build cannot resolve it.

\section{The contract suite (the behavioural spec)}

Nineteen contract documents, one per feature (001--015), specify behaviour at the
seam where it can be enforced:

\begin{center}
\begin{longtable}{@{}p{0.46\textwidth}p{0.48\textwidth}@{}}
\toprule
\textbf{Contract} & \textbf{Specifies} \\
\midrule
\texttt{011 .../instruction-document.md} & the four-section shape the loader requires \\
\texttt{011 .../compliance-verdict.md} & the verdict and LOCAL/ACCUMULATED attribution \\
\texttt{011 .../stage-execution.md} & the stage seam and its failure policy \\
\texttt{012 .../verification-result.md} & the three-state verifier result \\
\texttt{012 .../verification-observability.md} & how a substituted result is marked \\
\texttt{013 .../cost-record.md, cost-report.md} & cost record and report shape \\
\texttt{014 .../edits.md, gate.md, skill-document.md} & the skill edit ops, the gate, the doc format \\
\texttt{015 .../diagnostic-record.md, contribution-measurement.md, evolution-round.md} & diagnostics, leave-one-out, the round \\
\bottomrule
\end{longtable}
\end{center}

Together they are the platform's own specification of itself: each feature ships a
\texttt{spec.md}, a \texttt{data-model.md} and its contracts before its code, and the
tests are written against those documents.

\section{The skill format (the optimisation spec)}

The skill document (\texttt{contracts/skill-document.md}) has a fixed structure ---
title, granularity, ``when to apply'', numbered rules, and a protected
\texttt{SLOW\_UPDATE} region. It is the shape a reflective optimiser may edit and the
boundary it may not cross.

\section{Where each spec is consumed}

\begin{center}
\begin{longtable}{@{}p{0.42\textwidth}p{0.52\textwidth}@{}}
\toprule
\textbf{Specification} & \textbf{Consumed by} \\
\midrule
Blueprint / spec.md & \texttt{spec\_service} (parse, validate, store), every stage's task payload \\
Constitution & compliance gate, stage instruction docs, repair cap \\
Instruction documents & \texttt{instructions.py} loader, \texttt{render\_stage\_request} \\
Dependency allowlist & scaffolder compliance check \\
Contracts & the test suite (written against them), data models \\
Skill document & \texttt{skills/document.py}, the injection prefix \\
\bottomrule
\end{longtable}
\end{center}

\end{document}
